% RESULTADO_RECUPERADO: reunión local de propietarios ya impresos en el
% capítulo 8. No introduce un axioma ni una versión reducida de APP.
\begin{theorem}[Continuité conjointe de l'holonomie nonadique]
\label{thm:nucleo-continuidad-conexion-nonadica-conjunta}
Soit \(\widehat X_K\) l'état APP--TRIT--TPK enrichi par feuillet,
orientation, résidu, quotient, retenue, signature, cylindre, frontière, chemin,
mémoire et survie. La composition des neuf relations locales
publie une unique holonomie
\[
 \Gamma_{9,K}=\mathcal R_{K+8}\circ\cdots\circ\mathcal R_K
 :\widehat X_K\longrightarrow\widehat X_{K+9},
\]
compatible avec les troncatures du système projectif. À chaque tour, la phase revient
et la mémoire avance sans réinitialiser l'état ; la retenue conserve
sa loi de cocycle, le transfert agit simultanément sur ses neuf
composantes et la réalisation opératorielle conserve tant la composante
orthogonale que le terme terminal de l'identité de Stein. Par composition,
ces propriétés déterminent des sections compatibles à profondeur arbitraire,
non des choix indépendants niveau par niveau.
\end{theorem}

\begin{proof}
La composition et le retour sans réinitialisation sont, respectivement,
\cref{p12-83-c7-s1:def:holonomia-nonadica,p12-83-c7-s1:thm:retorno-nonadico-sin-reinicio}.
L'action conjointe sur les neuf coordonnées et leur transfert sont
\cref{p12-83-c7-s1:thm:funciones-nueve-componentes,p12-83-c7-s1:prop:transferencia-conjunta-nueve}.
La décomposition compression--expression et la conservation du terminal sont
\cref{p12-83-c7-s1:thm:conexion-compresion-expresion,p12-83-c7-s1:thm:stein-telescopia-terminal}.
La prolongation coinductive certifiée dans ce même chapitre conserve ces
identités sous troncature ; la famille compatible appartient donc à la
limite projective de l'état enrichi et peut être itérée à toute profondeur.
\end{proof}
